\subsubsection{Sprint 10}

\subsubsubsection{Informazioni generali e attività da svolgere}

\begin{tabular}{p{0.25\textwidth} p{0.2\textwidth}}
	\textbf{Inizio} & 16/03/2026 \\
	\textbf{Fine prevista} & 23/03/2026 \\
	\textbf{Fine reale} & 23/03/2026 \\
	\textbf{Giorni di ritardo} & 0
\end{tabular}
\newline \newline 
Le nostre attività da svolgere, definite nel \href{https://nullpointersgroup.github.io/Documentazione/output/PB/Verbali%20Interni/2026-03-16_verbale_interno.pdf}{Verbale interno del 16/03/2026} sono:
\begin{itemize}[itemsep=5pt, parsep=1pt, label=$\scriptstyle\bullet$]
	\item Scrivere il consuntivo$^G$ dello sprint$^G$ 9 e il preventivo dello sprint$^G$ 10
	\item Stesura Diario di Bordo del 18 marzo 2026
	\item Stesura schema del DB$^G$
	\item Stesura autenticazione e registrazione in Specifica tecnica$^G$
	\item Implementazione registrazione e autenticazione
	\item Inizio implementazione chat
\end{itemize}

\subsubsubsection{Rischi attesi}
I possibili rischi attesi sono:
\begin{itemize}[itemsep=5pt, parsep=1pt, label=$\scriptstyle\bullet$]
	\item RP3: rischio d'incomprensione di qualche task$^G$
	\item RA1: rischio di qualche imprevisto personale
	\item RE1: rischio dovuto alla poca conoscenza sulle nuove tecnologie da affrontare
	\item RE2: rischio dovuto alla stima errata delle attività
\end{itemize}

\subsubsubsection{Preventivo}
Si prospetta l'utilizzo delle seguenti risorse:
\begin{table}[H]
	\centering
	\makebox[\textwidth][c]{%
		\begin{minipage}{1.2\textwidth}
			\centering
			% Prima tabella senza bordi per le intestazioni ruotate
			\begin{tabular}{p{0.2\textwidth}p{0.12\textwidth}p{0.12\textwidth}p{0.12\textwidth}p{0.12\textwidth}p{0.12\textwidth}p{0.12\textwidth}}
				& \rotatebox{45}{\textbf{Responsabile$^G$}} &
				\rotatebox{45}{\textbf{Amministratore$^G$}} & \rotatebox{45}{\textbf{Analista$^G$}} & \rotatebox{45}{\textbf{Programmatore$^G$}} & \rotatebox{45}{\textbf{Verificatore$^G$}} &
				\rotatebox{45}{\textbf{Progettista$^G$}} \\
			\end{tabular}
			
			\vspace{0.2cm} % Piccolo spazio tra le due tabelle
			
			% Seconda tabella con i dati e i bordi
			\begin{tabular}{|p{0.2\textwidth}|p{0.12\textwidth}|p{0.12\textwidth}|p{0.12\textwidth}|p{0.12\textwidth}|p{0.12\textwidth}|p{0.12\textwidth}|}
				\hline
				M. Mazzaretto & - & - & - & 4 & 3 & 5 \\ \hline
				T. Ceron      & - & - & - & 4 & 3 & 5 \\ \hline
				L. Pieripolli & - & - & - & 4 & 3 & 5 \\ \hline
				L. Marcuzzo   & 4 & - & - & - & 1 & - \\ \hline
				M. Brunello   & - & - & - & 4 & 3 & 5 \\ \hline
				L. Casagrande & - & - & - & 4 & 3 & 5 \\ \hline
			\end{tabular}
		\end{minipage}
	}
	\caption{Sprint 10: Preventivo}
\end{table}

Il totale preventivato per lo sprint$^G$ è \textbf{1285€}.\\

\begin{figure}[H]
	\centering
	\includegraphics[width=0.8\textwidth]{PianoProgetto/sprint10_preventivo}
	\caption{Sprint 10: Preventivo}
\end{figure}

\subsubsubsection{Consuntivo}
In questo sprint$^G$ sono state utilizzate le seguenti risorse:
\begin{table}[H]
	\centering
	\makebox[\textwidth][c]{%
		\begin{minipage}{1.2\textwidth}
			\centering
			% Prima tabella senza bordi per le intestazioni ruotate
			\begin{tabular}{p{0.2\textwidth}p{0.12\textwidth}p{0.12\textwidth}p{0.12\textwidth}p{0.12\textwidth}p{0.12\textwidth}p{0.12\textwidth}}
				& \rotatebox{45}{\textbf{Responsabile$^G$}} &
				\rotatebox{45}{\textbf{Amministratore$^G$}} & \rotatebox{45}{\textbf{Analista$^G$}} & \rotatebox{45}{\textbf{Programmatore$^G$}} & \rotatebox{45}{\textbf{Verificatore$^G$}} &
				\rotatebox{45}{\textbf{Progettista$^G$}} \\
			\end{tabular}
			
			\vspace{0.2cm} % Piccolo spazio tra le due tabelle
			
			% Seconda tabella con i dati e i bordi
			\begin{tabular}{|p{0.2\textwidth}|p{0.12\textwidth}|p{0.12\textwidth}|p{0.13\textwidth}|p{0.12\textwidth}|p{0.12\textwidth}|p{0.12\textwidth}|}
				\hline
				M. Mazzaretto & - & - & - & 8 \textcolor{red}{(+4)}& 4 \textcolor{red}{(+1)} & 5 \\ \hline
				T. Ceron      & - & - & - & 4  & 2 \textcolor{green}{(-1)} & 8 \textcolor{red}{(+3)} \\ \hline
				L. Pieripolli & - & - & - & 6 \textcolor{red}{(+2)}& 3 & 4 \textcolor{green}{(-1)} \\ \hline
				L. Marcuzzo   & 3 \textcolor{green}{(-1)} & - & - & - & 1 & - \\ \hline
				M. Brunello   & - & - & - & 4 & 1 \textcolor{green}{(-2)} & 8 \textcolor{red}{(+3)} \\ \hline
				L. Casagrande & - & - & - & 0 \textcolor{green}{(-4)} & 3 & 8 \textcolor{red}{(+3)} \\ \hline
			\end{tabular}
		\end{minipage}
	}
	\caption{Sprint 10: Consuntivo}
\end{table}
Il totale consuntivato per lo sprint$^G$ è \textbf{1455€}.\\

In questo sprint$^G$, il carico di lavoro effettivo ha superato le aspettative in termini di ore produttive, poiché il gruppo ha dovuto dedicare maggiore impegno alle attività, in quanto il gruppo ha dovuto riprogettare e ricodificare parti ingenti del sistema in quanto il dal ricevimento con il prof. Cardin sono emerse gravi problematiche architetturali.
\begin{figure}[H]
	\centering
	\includegraphics[width=0.8\textwidth]{PianoProgetto/sprint10_consuntivo}
	\caption{Sprint 10: Consuntivo}
\end{figure}

\subsubsubsection{Retrospettiva}
Il gruppo, nonostante sia "tornato indietro" nella progettazione, è riuscito comunque a completare le attività che si era preffisato.
\subsubsubsection{Rischi incontrati}
In questo sprint$^G$ è stato riscontrato il rischio RE1 in quanto il gruppo ha dovuto interrompere l'avanzamento per rimediare agli errori fatti precedentemente durante la fase di progettazione architetturale.

\subsubsubsection{Efficacia delle strategie di gestione di mitigazione dei rischi}
Il gruppo è riuscito ad arginare il danno dovuto al blocco degli avanzamenti: applicando i suggerimenti forniti dal prof. Cardin durante il ricevimento. 
Permettendo così di completare le attività prefissate per la fine dello sprint.

\subsubsubsection{Migliorie da attuare per le attività future}
Nei prossimi sprint il gruppo dovrà prestare più attenzione in fase di progettazione, in particolar modo stando attenti nell'applicare i design patterns qualora si presenti l'effettivo problema che essi mirano a risolvere, senza forzare il pattern dove non è necessario.

\subsubsubsection{Aggiornamento delle risorse rimanenti}
\begin{table}[H]
	\centering
	\begin{tabular}{|p{0.2\textwidth}|p{0.1\textwidth}|p{0.05\textwidth}|p{0.1\textwidth}|p{0.15\textwidth}|p{0.17\textwidth}|}
		\hline
		\rowcolor{gray!25}
		Ruolo & Costo & Ore & Costo \newline effettivo & Ore \newline rimanenti & Budget \newline rimanente \\ \hline
		
		Responsabile$^G$
		& 30(€/h) & 3 & 90
		& 20 \textcolor{red}{(-3)}
		& 600 \textcolor{red}{(-90)} \\ \hline
		
		Amministratore$^G$
		& 20(€/h) & - & -
		& 1
		& 20 \\ \hline
		
		Analista$^G$
		& 25(€/h) & - & -
		& 2.5
		& 62.5\\ \hline
		
		Progettista$^G$
		& 25(€/h) & 33 & 825
		& 69 \textcolor{red}{(-33)}
		& 1725 \textcolor{red}{(-825)}\\ \hline
		
		Verificatore$^G$
		& 15(€/h) & 14 & 210
		& 73 \textcolor{red}{(-14)}
		& 1095 \textcolor{red}{(-210)} \\ \hline
		
		Programmatore$^G$
		& 15(€/h) & 22 & 330
		& 75 \textcolor{red}{(-22)}
		& 1125 \textcolor{red}{(-330)} \\ \hline
		
		Totale
		& - & 42 & 1455
		& 240.5 \textcolor{red}{(-72)}
		& 4627.5 \textcolor{red}{(-1455)} \\ \hline
	\end{tabular}
	\caption{Sprint 10: Aggiornamento risorse}
\end{table}